\uuid{x9YC}
\titre{Continuité en un point}
\auteur{bodin ; hanani ; mzari}
\organisation{exo7}
\licence{CC BY-NC-SA 4.0}
\datecreate{2020-01-20}
\niveau{L1}
\module{Analyse}
\chapitre{Continuité, limite et étude de fonctions réelles}
\sousChapitre{Continuité : théorie}
\difficulte{3}
\theme{Fonctions continues, continuité en un point}
\exo7qcmid{lille-1.continuite.018}

\contenu{
\begin{question}
\qclassification{123.01,123.02}
\qtags{motcle=continuité en un point}
Soit $f : \Rr \to \Rr$ une fonction et $(u_n)_{n\in\Nn}$ une suite. Quelles sont les assertions vraies ?
\begin{answers}
    \good{Si $u_n \to \ell$ et $f$ continue en $\ell$, alors $f(u_n)$ admet une limite.}

    \bad{Si $f(u_n) \to f(\ell)$ et $f$ est continue en $\ell$, alors $u_n \to \ell$.}

    \good{Si $u_n \to \ell$ et $f(u_n)$ n'a pas de limite, alors $f$ n'est pas continue en $\ell$.}

    \good{Si pour toute suite qui vérifie $u_n \to \ell$, on a $f(u_n) \to f(\ell)$, alors $f$ est continue en $\ell$.}

\end{answers}
\begin{explanations}
Une fonction $f$ est continue en $\ell$ si et seulement si
pour toute suite $(u_n)$ qui tend vers $\ell$, on a $f(u_n) \to f(\ell)$.
\end{explanations}
\end{question}
}
